\chapter{Rules of calculus}
\label{chap:rules}

With the chain rule (Theorem~\ref{thm:chain-rule}) and the structure of the higher
differentials (Chapter~\ref{chap:higher}) at hand, we can now derive the remaining rules of
calculus that are used in practice. All of them are classical in the Banach setting, and their
statements are the familiar ones; what has to be checked is that the proofs survive without
norms. In each case the argument either reduces to the chain rule or, through a continuous
linear functional, to the calculus of real functions of one variable.

We begin with continuous bilinear maps and the product rule. We then compute the derivatives
of a map along a line, which leads to Taylor's formula in a weak form, valid without any
completeness assumption on $F$. The chapter closes with the chain rule for second
differentials. The references are the results of \cite{GN26} quoted in each proof.

\section{Bilinear maps and the product rule}

\begin{theorem}[Bilinear maps]
  \label{thm:bilinear}
  \lean{LocallyConvex.Analysis.hasDirDerivAt_bilin, LocallyConvex.Analysis.contDiffMBOn_bilin}
  \leanok
  \uses{def:ck, thm:partials}
  Let $\beta \colon A \times B \to F$ be continuous and bilinear. Then $\beta$ is $\Ck{k}$ for
  every $k$, with $\dd \beta\bigl((a, b), (h, w)\bigr) = \beta(a, w) + \beta(h, b)$.
\end{theorem}

\begin{proof}
  \leanok
  The difference quotient of $\beta$ at $(a,b)$ in the direction $(h,w)$ is
  $\beta(a, w) + \beta(h, b) + t\,\beta(h, w)$, and letting $t \to 0$ gives the formula. The
  differential is therefore a sum of continuous bilinear maps composed with continuous linear
  maps, and induction on $k$ completes the proof. This is \cite[1.2.3 and 1.3.5]{GN26}.
\end{proof}

Note that continuity of $\beta$ means joint continuity on $A \times B$. This cannot be
weakened to separate continuity: for $b = 0$ and $h = 0$ the differential above reduces to
$\beta(a, w)$, which must be jointly continuous in $(a, w)$.

\begin{corollary}[Product rule]
  \label{cor:product}
  \lean{LocallyConvex.Analysis.ContDiffMBOn.bilin, LocallyConvex.Analysis.ContDiffMBOn.dMB_bilin}
  \leanok
  \uses{thm:bilinear, thm:chain-rule, lem:prod}
  If $f \colon U \to A$ and $g \colon U \to B$ are $\Ck{k}$ and $\beta$ is continuous and
  bilinear, then $x \mapsto \beta(f(x), g(x))$ is $\Ck{k}$ and
  \[
    \dd\bigl(\beta(f, g)\bigr)(x)(h) = \beta\bigl(f(x), \dd g(x)(h)\bigr)
      + \beta\bigl(\dd f(x)(h), g(x)\bigr).
  \]
\end{corollary}

\begin{proof}
  \leanok
  It is immediate from the chain rule applied to $\beta \circ (f, g)$, together with
  Theorem~\ref{thm:bilinear}.
\end{proof}

\begin{example}
  Let $\beta \colon E \times E \to F$ be continuous and bilinear, and let
  $q(x) = \beta(x, x)$. Taking $f = g = \mathrm{id}_E$ in Corollary~\ref{cor:product}, we
  obtain that $q$ is smooth, with $\dd q(x)(h) = \beta(x, h) + \beta(h, x)$. When $\beta$ is
  symmetric this is the familiar $2\beta(x, h)$.
\end{example}

\section{Taylor's formula}

To pass from the differentials of $f$ to the values of $f$ we restrict $f$ to a line and
compose with a continuous linear functional. The resulting real function of one variable is
as regular as $f$, and its derivatives are the iterated differentials of $f$ along the line.

\begin{lemma}[Derivatives along a line]
  \label{lem:line}
  \lean{LocallyConvex.Analysis.ContDiffMBOn.hasDerivAt_iteratedDMB_line, LocallyConvex.Analysis.ContDiffMBOn.continuousOn_iteratedDMB_line}
  \leanok
  \uses{thm:iterated-regular}
  Let $f$ be $\Ck{k}$ on $U$, $j < k$ and $\psi$ a continuous linear functional on $F$. Then
  the real function $s \mapsto \psi\bigl(\dd^{(j)} f(x + sh; h, \dots, h)\bigr)$ has derivative
  $\psi\bigl(\dd^{(j+1)} f(x + sh; h, \dots, h)\bigr)$ wherever $x + sh \in U$.
\end{lemma}

\begin{proof}
  \leanok
  By the chain rule along the line $s \mapsto (x + sh, h, \dots, h)$, whose derivative is
  $(h, 0, \dots, 0)$.
\end{proof}

\begin{theorem}[Taylor's formula, weak form]
  \label{thm:taylor}
  \lean{LocallyConvex.Analysis.ContDiffMBOn.taylor}
  \leanok
  \uses{lem:line}
  Let $f$ be $\Ck{k+1}$ on a set containing the segment $[x, x+h]$. For every continuous linear
  functional $\psi$ on $F$,
  \[
    \psi\bigl(f(x + h)\bigr) = \sum_{j=0}^{k} \frac{1}{j!}\,
      \psi\bigl(\dd^{(j)} f(x; h, \dots, h)\bigr)
      + \int_0^1 \frac{(1-t)^k}{k!}\,
      \psi\bigl(\dd^{(k+1)} f(x + th; h, \dots, h)\bigr)\, dt .
  \]
\end{theorem}

\begin{proof}
  \leanok
  We apply Taylor's formula with integral remainder to the real function
  $s \mapsto \psi(f(x + sh))$, whose derivatives are given by Lemma~\ref{lem:line}. Compare
  \cite[1.6.1 and 1.6.14]{GN26}, where the formula is stated for the integrals in $F$; the weak form
  needs no completeness of $F$.
\end{proof}

\begin{remark}
  The remainder in Theorem~\ref{thm:taylor} is a real integral, one for each functional $\psi$.
  Writing it as a vector of $F$ requires the weak integral of
  $t \mapsto \frac{(1-t)^k}{k!}\,\dd^{(k+1)} f(x + th; h, \dots, h)$ to exist, which in general
  depends on a completeness property of $F$ (Chapters~\ref{chap:integrals}
  and~\ref{chap:completeness}). Since continuous linear functionals separate points, the weak
  form determines $f(x + h)$ all the same.
\end{remark}

\section{The second-order chain rule}

\begin{theorem}[Second-order chain rule]
  \label{thm:chain-second}
  \lean{LocallyConvex.Analysis.dMB_dMB_comp}
  \leanok
  \uses{thm:chain-rule}
  Let $\tau \colon U \to V$ and $g \colon V \to F$ be $\Ck{2}$, with $U \subseteq E$ and
  $V \subseteq E'$ open. For $y \in U$ and $a, b \in E$,
  \[
    \dd^{(2)} (g \circ \tau)(y; a, b)
      = \dd^{(2)} g\bigl(\tau y; \dd\tau(y)(a), \dd\tau(y)(b)\bigr)
        + \dd g(\tau y)\bigl(\dd^{(2)} \tau(y; a, b)\bigr).
  \]
\end{theorem}

The first term is what one would guess by applying the chain rule twice naively; the second
involves the second differential of $\tau$, and vanishes when $\tau$ is the restriction of a continuous
linear map, since then $\dd\tau(y)(b)$ does not depend on $y$.

\begin{proof}
  \leanok
  On $U \times E$, the first-order chain rule gives $\dd(g \circ \tau) = \dd g \circ \Phi$ with
  $\Phi(y', b') = (\tau y', \dd\tau(y')(b'))$. We differentiate this once more with the
  first-order chain rule. The derivative of $\Phi$ at $(y, b)$ in the direction $(a, 0)$ is
  $(\dd\tau(y)(a), \dd^{(2)}\tau(y; a, b))$. Splitting it as
  $(\dd\tau(y)(a), 0) + (0, \dd^{(2)}\tau(y; a, b))$, the first term gives $\dd^{(2)} g$, and
  the second gives $\dd g$, because $\dd g(\tau y)$ is linear.
\end{proof}
